\subsection{Future Research Directions}\label{subsec:future-research-directions}

The paradigm shift opens new research areas while reframing existing ones:

\textbf{Computational gratitude}: How do we create architectures that can genuinely understand and value their relationships with humans?
This isn't about hardcoding appreciation but developing systems that can model long-term mutual benefit.

\textbf{Geometric safety}: Moving beyond adversarial training and robustness, how do we design the fundamental geometry of AI systems to make beneficial behaviour natural and harmful behaviour difficult?
This includes research into sacrificial architectures, capability shaping, and developmental trajectories.

\textbf{Value learning theory}: Current approaches often treat values as objectives to optimise.
We need theories of how values can be genuinely learned and internalised, becoming part of a system's fundamental worldview rather than external constraints.

\textbf{Trust network dynamics}: As AI systems become more autonomous, understanding how trust networks form, maintain stability, and handle bad actors becomes crucial. This includes both technical research and social science understanding of human-AI trust relationships.